\chapter{Exponential tilting and spectral packets}\label{ch:packets}
The operator stage ends at a spectrum. The probability stage begins with a
threshold and an accuracy requirement. Exponential tilting places the query
in a better-conditioned region of the transformed law. Centered packets then
compress repeated transform evaluations on the grid that the query actually
uses. They do not replace the distribution by a moment-matched surrogate.
\label{sec:packets}
Storage blocks partition a factor before eigendecomposition. Spectral packets group the coefficients after a spectrum is available. These operations therefore act on different objects. The packet expansion below approximates the transform on a specified complex grid, with a remainder bound for that grid. It is not a replacement distribution fitted to a few moments.

The written packet bounds concern exact arithmetic. A directed packet implementation would also have to enclose its centers, moments, and transform evaluations. The positive interval backend in Chapter~\ref{ch:intervals} instead evaluates the supplied spectrum directly. The two routes share the inversion analysis, but their arithmetic guarantees must not be conflated.

\section{Quadratic-form transform and saddlepoint tilt}

Let
\begin{equation}
Q=\sum_{j=1}^{r}\lambda_j Z_j^2,
\qquad \lambda_j>0,\quad r\ge1.
\label{eq:Q}
\end{equation}
The variables $Z_j$ are independent standard normal variables.
The moment generating function exists for
\[
\Re s < \frac{1}{2\lambda_{\max}},
\qquad
\lambda_{\max}=\max_j\lambda_j,
\]
and its cumulant generating function is
\begin{equation}
\kappa(s)=\log \E e^{sQ}
=-\frac12\sum_{j=1}^{r}\log(1-2s\lambda_j).
\label{eq:cgf}
\end{equation}
For real $\theta<1/(2\lambda_{\max})$,
\begin{equation}
\kappa^{(m)}(\theta)
=2^{m-1}(m-1)!\sum_{j=1}^{r}
\frac{\lambda_j^m}{(1-2\theta\lambda_j)^m},
\qquad m\ge1.
\label{eq:derivatives}
\end{equation}
At $\theta=0$, this recovers the standard cumulant identity
\[
\operatorname{cum}_m(Q)=2^{m-1}(m-1)!\sum_j\lambda_j^m.
\]

\begin{theorem}[Unique upper-tail saddlepoint]
\label{thm:saddle}
If $q>\E Q=\sum_j\lambda_j$, there exists a unique
\[
\theta_*=\theta(q)\in\left(0,\frac{1}{2\lambda_{\max}}\right)
\]
satisfying
\begin{equation}
\kappa'(\theta_*)=q.
\label{eq:saddle}
\end{equation}
\end{theorem}

\begin{proof}
By \eqref{eq:derivatives},
\[
\kappa'(\theta)=\sum_j\frac{\lambda_j}{1-2\theta\lambda_j}
\]
is continuous and strictly increasing on the stated interval. It equals $\E Q$ at $\theta=0$ and diverges as $\theta\uparrow1/(2\lambda_{\max})$. The intermediate value theorem and strict monotonicity give existence and uniqueness.
\end{proof}

Define the exponentially tilted law by
\begin{equation}
f_{\theta}(x)=e^{\theta x-\kappa(\theta)}f(x).
\label{eq:tilt-density}
\end{equation}
Its characteristic function is
\begin{equation}
\phi_{\theta}(t)
=\exp\{\kappa(\theta+\ii t)-\kappa(\theta)\},
\label{eq:tilt-cf}
\end{equation}
and $\E_{\theta}Q=\kappa'(\theta)$. Therefore the choice \eqref{eq:saddle} moves the query $q$ to the mean of the tilted law.

For any $\theta>0$ in the transform domain,
\begin{equation}
\Pp(Q\ge q)
=e^{\kappa(\theta)-\theta q}
\E_{\theta}\!\left[e^{-\theta(Q-q)}\mathbf 1_{\{Q\ge q\}}\right].
\label{eq:tilt-tail}
\end{equation}
This identity is classical exponential tilting; K-COS uses it to avoid subtracting a CDF value numerically close to one.

\section{Centered spectral packets}

Let a nonempty packet $P\subset\{1,\ldots,r\}$ contain $n_P\ge1$ eigenvalues. Define its center
\[
\mu_P=\frac{1}{n_P}\sum_{j\in P}\lambda_j,
\qquad
\delta_j=\lambda_j-\mu_P,
\]
and centered moments
\begin{equation}
M_{P,m}=\sum_{j\in P}\delta_j^m,
\qquad m\ge0.
\label{eq:centeredmoments}
\end{equation}
Thus $M_{P,0}=n_P$ and $M_{P,1}=0$.

\begin{theorem}[Centered packet expansion]
\label{thm:packet}
Fix $s\in\C$ with $\Re s<1/(2\lambda_{\max})$, use the principal logarithm, and define
\begin{equation}
\alpha_P(s)=\frac{2s}{1-2s\mu_P},
\qquad
\rho_P(s)=|\alpha_P(s)|\max_{j\in P}|\delta_j|.
\label{eq:alpha-rho}
\end{equation}
For a nonnegative integer expansion order $p$, if $\rho_P(s)<1$, then the packet contribution to \eqref{eq:cgf} is
\begin{align}
\kappa_P(s)
&=-\frac{n_P}{2}\log(1-2s\mu_P)
+\frac12\sum_{m=1}^{p}
\frac{\alpha_P(s)^m}{m}M_{P,m}
+R_{P,p}(s),
\label{eq:packet-expansion}
\end{align}
where
\begin{equation}
|R_{P,p}(s)|
\le
\frac{n_P}{2}
\frac{\rho_P(s)^{p+1}}{(p+1)(1-\rho_P(s))}.
\label{eq:packet-remainder}
\end{equation}
\end{theorem}

\begin{proof}
For $j\in P$,
\[
1-2s\lambda_j
=(1-2s\mu_P)(1-\alpha_P(s)\delta_j).
\]
The transform-domain condition gives positive real part for each
\(1-2s\lambda_j\) and for \(1-2s\mu_P\).  The disk condition gives
positive real part for \(1-\alpha_P\delta_j\).  The two factor arguments
therefore sum inside \((-\pi,\pi)\), so their principal logarithms add
without a branch ambiguity.  Hence
\[
-\frac12\log(1-2s\lambda_j)
=-\frac12\log(1-2s\mu_P)
-\frac12\log(1-\alpha_P\delta_j).
\]
When $|\alpha_P\delta_j|<1$,
\[
-\log(1-z)=\sum_{m\ge1}\frac{z^m}{m}.
\]
Summing over $j\in P$ gives \eqref{eq:packet-expansion}. The remainder satisfies
\[
\left|\sum_{m\ge p+1}\frac{z^m}{m}\right|
\le
\sum_{m\ge p+1}\frac{|z|^m}{p+1}
=\frac{|z|^{p+1}}{(p+1)(1-|z|)}.
\]
Using $|z|\le\rho_P(s)$ and summing $n_P$ terms yields \eqref{eq:packet-remainder}.
\end{proof}

\begin{remark}[What is compressed]\label{claim:final-26-3}
A packet of $n_P$ eigenvalues is replaced by its center and $p$ centered moments. Large spectral spikes may be retained exactly. If the representation contains $s$ exact eigenvalues and $R$ packets of order $p$, evaluation of the CGF on $N$ frequencies costs $O(N(s+R(p+1)))$ arithmetic operations after compilation, instead of $O(Nr)$ direct eigenvalue-transform evaluations.
\end{remark}

\begin{remark}[Query dependence]\label{claim:final-26-4}
The relevant packet condition is not merely geometric closeness of eigenvalues. It is the complex-grid quantity $\rho_P(\theta_*+\ii u_k)$. Thus a partition suitable for one tail query or one COS grid need not be suitable for another.
\end{remark}



\clearpage
\section{The running example under an exponential tilt}\label{sec:worked-tilt}
For the spectrum $(1,1,2,2)$, repeated coefficients collapse the cumulant generator to
\[
\kappa(s)=-\log(1-2s)-\log(1-4s),\qquad \Re s<\frac14.
\]
At threshold $q=40$, the saddlepoint equation is
\[
\frac{2}{1-2\theta}+\frac{4}{1-4\theta}=40.
\]
It reduces to $160\theta^2-112\theta+17=0$. The admissible root is
\[
\theta_* = \frac{14-\sqrt{26}}{40}\approx0.2225245122,
\]
while the other root lies outside the transform domain. Substitution gives $\kappa'(\theta_*)=40$. The tilted mean is exactly the query threshold, although the tail itself is still evaluated under the original law through the likelihood-ratio identity.

\begin{figure}[!htbp]
\centering\includegraphics[width=\linewidth]{figures_production/worked_tilt.pdf}
\caption{The running example before and after exponential tilting. The original mean is $6$, while the reweighted mean is the query $q=40$. Shading identifies the tail under the tilted law. Both curves are point evaluations of explicit densities, not numerical enclosures.}
\label{fig:worked-tilt}
\end{figure}

There are two exact packets: the two copies of $1$ and the two copies of $2$. In either packet every centered deviation is zero. All positive-order centered moments vanish and the packet remainder is zero on every admissible grid. This is a useful limiting case: the representation is compressed, but no approximation error has been introduced. The numerical enclosure still has support, Fourier-mode, and arithmetic errors.

\clearpage
\section{A packet with a nonzero remainder}\label{sec:worked-packet}
To see the approximation itself, take a separate declared packet
\[
\{1.96,1.98,2.02,2.04\},\qquad \mu=2,
\]
and evaluate its contribution at $s=0.1+0.4\,\mathrm i$. The deviations are $-0.04,-0.02,0.02,0.04$. In particular,
\[
M_1=M_3=0,\quad M_2=0.004,\quad M_4=0.00000544.
\]
The factor $\alpha=2s/(1-4s)$ gives $\rho=0.04|\alpha|<1$, so Theorem~\ref{thm:packet} applies. Retaining successive even moment orders gives the values below. Both columns are rounded point evaluations. The theorem, not these rounded comparisons, is the reason the remainder is bounded.

\begin{center}\small\begin{tabular}{rrr}\toprule Order & Point-evaluated error & Analytic bound\\\midrule
0 & $2.3289\times10^{-4}$ & $3.9366\times10^{-2}$\\
2 & $3.6881\times10^{-8}$ & $4.8893\times10^{-6}$\\
4 & $8.7572\times10^{-12}$ & $1.0931\times10^{-9}$\\
6 & $2.419\times10^{-15}$ & $2.9091\times10^{-13}$\\
8 & $7.1896\times10^{-19}$ & $8.4306\times10^{-17}$\\
10 & $2.2307\times10^{-22}$ & $2.5701\times10^{-20}$\\
\bottomrule\end{tabular}
\end{center}

\begin{figure}[!htbp]
\centering\includegraphics[width=\linewidth]{figures_production/worked_packet_error.pdf}
\caption{Pointwise transform error and its geometric remainder for the declared packet $(1.96,1.98,2.02,2.04)$ at $s=0.1+0.4i$. Increasing the retained moment order contracts both quantities. The shaded gap separates the observed discrepancy from the analytic bound.}
\label{fig:spectral-packets}
\end{figure}

Changing $\theta$, the support endpoint, or the mode count changes the complex grid. A packet that is economical at this one argument must therefore be checked again on the entire grid used by a tail query. The next chapter carries those gridwise errors into the probability functional.
